%% F04 --- Sums, products and sequences
%%
%% Every number in this program that the reader cannot do in their head comes
%% from code/f04_sums.py and is pulled in with \val{}. Arithmetic the program
%% is TEACHING -- 2 + 5 + 1 + 7, 4! = 24, 8 x 9 / 2 -- is written inline,
%% because the reader is meant to do it.
%%
%% THE DEBT THIS PROGRAM PAYS. Program F03 wrote the log-probability of a
%% sequence as a sigma over its tokens and then divided by the token count,
%% and it never defined that sigma; F02 hands it over by name -- "Program F4
%% reads a sum as a loop" -- and F02 also introduces the subscript. Section 4
%% opens by naming that debt and quoting F03's own values rather than new
%% ones, so the three programs are provably quoting one computation.
%%
%% TWO PAIRS OF KEYS THAT LOOK LIKE ONE NUMBER PRINTED TWICE AND ARE NOT.
%% f04.geom.ten and f04.ema.w10 differ by exactly (1 - beta); f04.geom.inf is
%% the same 10 as the EMA's effective window. Both are the program's punchline
%% arriving twice, and the frames say so rather than leaving the reader to
%% spot the coincidence. See the module docstring of code/f04_sums.py.
%%
%% Twin: programs/pl/F04-sums-products-sequences.tex. Frame for frame, macro
%% for macro, number for number; tools/parity.py compares them.

\program{Sums, products and sequences}
\label{prog:F04}
\index{sigma notation!as a loop}
\index{product!pi notation}

\begin{outcomes}
\outcome{1--6}{read a sigma aloud, write a short one out in full, and turn a
  loop with an accumulator into a sigma and a sigma back into a loop}
\outcome{7--11}{say how many terms a sum has from its limits alone, re-index a
  sum so that it starts where you need it to, and count the query--key pairs a
  causal mask computes}
\outcome{12--16}{write a repeated multiplication in pi notation, and say what
  an empty sum and an empty product each come to and why the answer is forced
  rather than chosen}
\outcome{17--22}{write mean squared error and cross-entropy in the same shape,
  and point at the sigma and the divisor in any loss you are shown}
\outcome{23--28}{work out an average over items from a set of group averages
  and the sizes of their groups, and say which of the two averages a reported
  number is}
\outcome{29--34}{write a sequence as a recurrence and as a general formula, sum
  a geometric series, and use a logarithm to turn a long product into a sum}
\outcome{35--42}{write an exponential moving average as a two-line recurrence
  and as a weighted sum, say what its coefficient controls, and work out what
  it reads after $t$ steps from a given starting value}
\end{outcomes}

\begin{quiz}
\item Write $\sum_{i=1}^{4} 2i$ out in full and evaluate it. \teachesat{1--2}
  \answerto{$2 + 4 + 6 + 8 = 20$. The body is an expression in $i$, evaluated
  once for each value the index takes.}
\item How many terms are there in $\sum_{i=0}^{n} x_i$? \teachesat{7--8}
  \answerto{$n + 1$. Both limits are included, so the count is
  $\text{stop} - \text{start} + 1$.}
\item Rewrite $\sum_{i=1}^{n} x_i$ so that the index starts at $0$.
  \teachesat{9--10}
  \answerto{$\sum_{j=0}^{n-1} x_{j+1}$. Shift the limits down by one and the
  subscript up by one, so that the same terms are named.}
\item Is $\sum_{i=1}^{n} a_i b_i$ equal to
  $\left(\sum_{i=1}^{n} a_i\right)\left(\sum_{i=1}^{n} b_i\right)$?
  \teachesat{5--6}
  \answerto{No. At $a = (1, 2)$ and $b = (3, 4)$ the left side is $11$ and the
  right side is $21$.}
\item What do \code{sum([])} and \code{math.prod([])} return? \teachesat{13--14}
  \answerto{$0$ and $1$. Each is the value that leaves the next term alone,
  which is the identity element of the operation.}
\item Write mean squared error over $n$ examples in sigma notation.
  \teachesat{17--18}
  \answerto{$\frac{1}{n}\sum_{i=1}^{n} (\hat{y}_i - y_i)^{2}$.}
\item A batch of $1000$ tokens has mean loss $\num{2.0}$ and a batch of $10$
  tokens has mean loss $\num{8.0}$. What is the mean loss over the $1010$
  tokens? \teachesat{23--24}
  \answerto{$\val{f04.batch.pooled}$. The two group means may not be averaged:
  one of them stands for a hundred times as many tokens as the other.}
\item Two evaluation shards score $\num{0.90}$ and $\num{0.10}$. Under what
  condition is the average of those two numbers the accuracy on the pooled
  items? \teachesat{27--28}
  \answerto{When the two shards hold the same number of items. Then, and only
  then, equal weights are the right weights.}
\item How many query--key pairs does a causal mask compute over $n$ positions?
  \teachesat{10--11}
  \answerto{$\frac{n(n+1)}{2}$, which is a little over half of $n^{2}$.}
\item Evaluate $\sum_{k=0}^{9} (\num{0.9})^{k}$, and say what it tends to as the
  number of terms grows. \teachesat{31--32}
  \answerto{$\val{f04.geom.ten}$, tending to $\val{f04.geom.inf}$.}
\item Write $\ln$ of a product of $n$ factors as a sum. \teachesat{33--34}
  \answerto{$\ln \prod_{i=1}^{n} x_i = \sum_{i=1}^{n} \ln x_i$, for positive
  factors.}
\item A running average is kept as $m \leftarrow \num{0.9}\,m + \num{0.1}\,g$.
  It currently reads $\num{1.0}$ and the new $g$ is $\num{11.0}$. What does it
  read after one step? \teachesat{35--36}
  \answerto{$\num{2.0}$. The coefficient $\num{0.9}$ weights the \emph{old}
  value, so nine-tenths of the history survives and one-tenth of the news gets
  in.}
\item Which forgets a change faster, $\beta = \num{0.9}$ or
  $\beta = \num{0.999}$? \teachesat{36--37}
  \answerto{$\beta = \num{0.9}$. The larger the coefficient, the more of the
  old value is kept and the slower the average moves.}
\item An exponential moving average starts at $0$ and every observation is
  exactly $1$. What does it read after one step at $\beta = \num{0.999}$?
  \teachesat{39--40}
  \answerto{$\num{0.001}$, a thousandth of the truth.}
\end{quiz}

% ============================================================
\section{A sum written short}
\label{sec:F04-sigma}
\index{sigma notation!four parts of}
\index{sigma notation!weighted sum}

\begin{fr}
Here is a loop you have written more times than you can remember. It totals a
list of per-example losses:

\begin{python}
total = 0.0
for x in losses:
    total = total + x
\end{python}

And here is the same instruction as a paper writes it:
\[ \sum_{i=1}^{n} x_i \]

That is not a new operation and there is nothing in it to be afraid of.
\textbf{A sigma is a loop written short.} The capital Greek sigma is doing the
work of the word \emph{for} and the word \emph{add} at the same time, and it
has four parts you can point at:

\begin{center}
\begin{tabular}{ll}
\toprule
Part & What it says \\
\midrule
$i$ & the name of the index, the thing that counts \\
$i = 1$, under the sign & where the count begins \\
$n$, above the sign & where it stops \\
$x_i$, to the right & the body, worked once per value of $i$ \\
\bottomrule
\end{tabular}
\end{center}

Said out loud it is \emph{the sum, for $i$ from $1$ to $n$, of $x$ sub $i$},
and a formula you can pronounce is one you can hold in your head.

The one thing the notation has that the loop has not is a name for the
\emph{position} of each term. Program~\ref{prog:F02} gave you the subscript
for exactly that \dash{} $x_i$ is the $i$th input \dash{} and here it is what
lets a formula talk about the third loss rather than only about all of them.

Write $\sum_{i=1}^{4} x_i$ out in full for $x = (2, 5, 1, 7)$, and give the
total.
\nextframe
\end{fr}

\begin{fr}
\ans{$15$}
Written out, $\sum_{i=1}^{4} x_i = 2 + 5 + 1 + 7 = 15$. Four values of the
index, four terms, one addition each, and no technique to acquire.

The body does not have to be a lookup. It is an expression in $i$, and it is
worked afresh for every value the index takes:
\[ \sum_{i=1}^{3} i^{2} \]
says put $i = 1$ into $i^{2}$, then $i = 2$, then $i = 3$, and add the three
results.

Fill in the three terms, and then the total:
\[ \sum_{i=1}^{3} i^{2} = \blank + \blank + \blank = \blank \]
\dotline
\nextframe
\end{fr}

\begin{fr}
\begin{ansblock}
$1 + 4 + 9 = 14$.
\end{ansblock}

Now drive the machine the other way, because that is the direction you need
when you are reading somebody else's code and want to say in one line what it
computes:

\begin{python}
total = 0.0
for wi, xi in zip(w, x):
    total = total + wi * xi
\end{python}

Write that loop as a sigma.
\nextframe
\end{fr}

\begin{fr}
\begin{ansblock}
\[ \sum_{i=1}^{n} w_i x_i \]
\end{ansblock}

It has a name \dash{} a \emph{weighted sum} \dash{} and it is one unit of a
linear layer, the dot product Program~\ref{prog:P05} is built on, and the
shape this book returns to more often than any other.

Two rules come with it, and neither is a new fact. A sigma splits over a sum:
\[ \sum_{i=1}^{n} (a_i + b_i) = \sum_{i=1}^{n} a_i + \sum_{i=1}^{n} b_i \]
and a constant factor comes out of the front:
\[ \sum_{i=1}^{n} c\,a_i = c \sum_{i=1}^{n} a_i \]

Both are the loop rearranged. The first says you may add each pair and then
total, or total each list and then add; the second says you may scale every
term or scale the answer. Check them on $a = (1, 2)$ and $b = (3, 4)$: the
first gives $4 + 6 = 10$ on the left and $3 + 7 = 10$ on the right, and with
$c = 5$ the second gives $5 + 10 = 15$ on the left and $5 \times 3 = 15$ on
the right.
\end{fr}

\begin{fr}
No working on this one and no checking. Write down the first answer that comes
to you, as fast as you can produce it.

\textbf{Is $\sum_{i=1}^{n} a_i b_i$ the same number as
$\left(\sum_{i=1}^{n} a_i\right)\left(\sum_{i=1}^{n} b_i\right)$?}

Yes or no.
\dotline
\nextframe
\end{fr}

\begin{fr}
\ans{No}
They are not even close, and the smallest example says so. Take $a = (1, 2)$
and $b = (3, 4)$, which is the pair you were holding a moment ago:
\[ \sum_{i=1}^{2} a_i b_i = 1 \times 3 + 2 \times 4 = 3 + 8 = 11 \]
\[ \left(\sum_{i=1}^{2} a_i\right)\left(\sum_{i=1}^{2} b_i\right)
   = 3 \times 7 = 21 \]

\begin{trapbox}
\textbf{Nothing distributed here, and that is the whole of it.} The two rules
of the previous frame are both about a \emph{sum} sitting inside the sigma:
split it, or pull a constant out of it. Neither says anything about a
\emph{product} sitting inside the sigma, and no rule does.

What produced the wrong answer is a true rule generalised past its hypothesis,
and it is the reasoning this book has caught before under other signs.
Program~\ref{prog:F01} has $\sqrt{a+b} \ne \sqrt{a} + \sqrt{b}$;
Program~\ref{prog:F02} has $(a+b)^{2} \ne a^{2} + b^{2}$; and
Program~\ref{prog:F03} has no rule at all for $\ln(a + b)$. Each of them is an
operation that distributes over a \emph{product} being asked to distribute over
a \emph{sum}.

So do not file this away as one more exception. File the question instead:
what is the operation, and what is it being applied to?
\end{trapbox}

\begin{aibox}
An attention score is $\sum_i q_i k_i$: a query against a key, one product per
coordinate, all of them added. It is not
$\left(\sum_i q_i\right)\left(\sum_i k_i\right)$, so there is no version of
attention in which you keep two running totals and multiply them at the end.
Every coordinate has to meet its partner, which is why a dot product costs a
multiplication per dimension and why Program~\ref{prog:P05} spends a whole
program on what that one number means.
\end{aibox}
\end{fr}

% ============================================================
\section{Where the index starts}
\label{sec:F04-index}
\index{sigma notation!re-indexing}
\index{troubleshooting!off-by-one in a sum}

\begin{fr}
The letter carries no meaning whatever. $\sum_{i=1}^{n} x_i$ and
$\sum_{j=1}^{n} x_j$ are the same number, and renaming the index changes
nothing, exactly as renaming a loop variable changes nothing. It is a bound
name: it exists between the two limits and nowhere else.

Where the index \emph{starts} is another matter, and the two disciplines you
work in disagree about it. Mathematics counts from one. Python counts from
zero. Here is one list of four numbers, labelled both ways:

\begin{center}
\begin{tabular}{lcccc}
\toprule
Value & $2$ & $5$ & $1$ & $7$ \\
\midrule
Mathematics & $x_1$ & $x_2$ & $x_3$ & $x_4$ \\
Python & \code{x[0]} & \code{x[1]} & \code{x[2]} & \code{x[3]} \\
\bottomrule
\end{tabular}
\end{center}

The same four values in the same order, under two sets of labels. Now a
question about the labels rather than about the values.

How many terms are there in $\sum_{i=0}^{n} x_i$?
\nextframe
\end{fr}

\begin{fr}
\ans{$n + 1$}

\begin{trapbox}
\textbf{This is the fence-post, and it is the oldest off-by-one there is.} A
fence a hundred metres long with a post every ten metres carries eleven posts,
not ten. Counting the gaps gives ten, and the fence has two ends.

A sigma has the same shape, because both limits are \emph{included}. The
number of terms is
\[ \text{stop} - \text{start} + 1 \]
and the answer $n$ is what you get by subtracting the limits and stopping
there. That subtraction counts the steps between the terms; the terms are one
more than the steps.
\end{trapbox}

Three shapes, and they are the three you will meet. Fill in the one that is
missing:

\begin{center}
\begin{tabular}{lll}
\toprule
Sum & Terms & The loop it is \\
\midrule
$\sum_{i=1}^{n}$ & $n$ & \code{range(1, n + 1)} \\
$\sum_{i=0}^{n-1}$ & \blank & \code{range(n)} \\
$\sum_{i=0}^{n}$ & $n + 1$ & \code{range(n + 1)} \\
\bottomrule
\end{tabular}
\end{center}

Then write down the sigma that \code{for i in range(5)} walks.
\dotline
\nextframe
\end{fr}

\begin{fr}
\begin{ansblock}
$n$ terms. And the loop walks
\[ \sum_{i=0}^{4} x_i \]
which is five terms, from $0$ to $4$ inclusive.
\end{ansblock}

\begin{notationbox}
\textbf{One English word doing two jobs.} Program~\ref{prog:F01} calls the
raised number in $a^{n}$ the \emph{index}, or the exponent. This program calls
the lowered number in $x_i$ the \emph{index} as well. The two are unrelated:
one says how many times to multiply, the other says which item you mean.

English lets them collide and leaves the reader to hold both. Polish does not:
the raised one is a \emph{wykładnik} and the lowered one an \emph{indeks}, two
words for two jobs. Knowing that the English word is overloaded is worth
something in either language, because the papers are written in English.
\end{notationbox}

Re-indexing is mechanical: move the limits one way and the subscript the other,
so that the same terms end up named. Written out in full, once:
\[ \sum_{i=1}^{n} x_i = \sum_{j=0}^{n-1} x_{j+1} \]
At $j = 0$ the term is $x_1$; at $j = n-1$ it is $x_n$. Same terms, same
order, same total.

Rewrite $\sum_{k=2}^{5} k^{2}$ so that the index starts at $0$.
\nextframe
\end{fr}

\begin{fr}
\begin{ansblock}
\[ \sum_{k=2}^{5} k^{2} = \sum_{j=0}^{3} (j+2)^{2} \]
Both sides come to $4 + 9 + 16 + 25 = 54$. The limits went down by two and the
subscript went up by two, which is the only move there is.
\end{ansblock}

This is not shuffling notation for its own sake. Positions in a tokeniser are
numbered from zero, positions in most papers are numbered from one, and a
model has both conventions running through it at once. Sitting between them is
the causal mask: row $i$ attends to the keys $0$ up to and including $i$, so
the first query sees one key, the second sees two, and the last sees all of
them.

Over $n = 8$ positions, how many query--key pairs does that mask compute?
\nextframe
\end{fr}

\begin{fr}
\begin{ansblock}
$36$, against $64$ for the full square:
\[ \sum_{i=1}^{8} i = 1 + 2 + \cdots + 8 = \frac{8 \times 9}{2} = 36 \]
\end{ansblock}

The closed form is Gauss's and the argument is one line. Write the sum
forwards and backwards underneath itself; every column adds to $n + 1$, there
are $n$ columns, and everything has been counted twice:
\[ \sum_{i=1}^{n} i = \frac{n(n+1)}{2} \]

At $n = 2048$ that is $\val{f04.causal.pairs}$ pairs against
$\val{f04.causal.sq}$ for the full square.

\begin{aibox}
\enquote{Causal masking halves the work} is an asymptotic statement quoted as
an exact one. The ratio is
\[ \frac{n(n+1)/2}{n^{2}} = \frac{n+1}{2n} \]
which falls towards a half from above and never arrives. At $n = 2048$ it is
$\val{f04.causal.ratio}$, near enough a half to say so; at $n = 8$ it is
$\val{f04.causal.ratio.eight}$, which is an eighth more work than half. The
excess is the diagonal \dash{} every query attends to its own key \dash{} so
it is a rounding error at long context and a real term at short.

Program~\ref{prog:P03} prices work in FLOPs rather than counting pairs, which
needs nothing beyond the arithmetic and the sigma you now have.
\end{aibox}
\end{fr}

\mermaidfig{f04-sigma-as-loop}{%
  The same instruction in two notations. The loop names no positions and the
  sigma does, which is why the sigma has to say where its index starts and the
  loop does not.}{%
  A loop, its four parts, and the two index conventions.}

% ============================================================
\section{Products, and the two empty cases}
\label{sec:F04-products}
\index{product!factorial}
\index{product!empty product}

\begin{fr}
The other symbol is the same machine with a different operation inside it.
Where the sigma adds, the capital Greek pi multiplies:

\begin{python}
total = 1.0
for x in values:
    total = total * x
\end{python}

\[ \prod_{i=1}^{n} x_i \]

Everything from the last two sections transfers untouched: the index and its
name, the two limits, the body, the fact that both limits are included, and
the business of re-indexing. Only the operation has changed.

Evaluate $\prod_{i=1}^{4} i$, and say what that number is usually called.
\nextframe
\end{fr}

\begin{fr}
\begin{ansblock}
$1 \times 2 \times 3 \times 4 = 24$, and it is $4!$, four factorial. A
factorial is a pi with the index itself in the body:
\[ n! = \prod_{i=1}^{n} i \]
\end{ansblock}

Now look back at the two loops. Each one starts its accumulator somewhere
before the first pass, and what it starts at is about to be worth a frame of
its own.

Quickly, and with no working. What do \code{sum([])} and \code{math.prod([])}
return?
\dotline
\nextframe
\end{fr}

\begin{fr}
\begin{ansblock}
$0$ and $1$.
\end{ansblock}

Not asserted \dash{} run:

\transcript{f04-empty}

The loop body never executes, so what comes back is whatever the accumulator
was set to, and the reason it was set to that is the same in both cases.
Adding $0$ leaves a number alone; multiplying by $1$ leaves a number alone.
Each is the \emph{identity element} of its operation, and it is the only value
that will do: the sum of no terms has to be the thing that leaves the next
term unchanged when it arrives, or the loop would give a different answer
depending on where you started it.

So what is $0!$?
\nextframe
\end{fr}

\begin{fr}
\ans{$0! = 1$}
It is the empty product, and it needs no argument of its own.

\begin{trapbox}
\textbf{\enquote{There is nothing there, so it is zero.}} Or, from the other
side, \enquote{there is nothing there, so it is undefined.} The first is
correct for one of the two operations and wrong for the other; the second is
wrong for both.

What produces it is that \emph{nothing} has one name and the two operations
have two identities. Emptiness is not a value, so it cannot have a value of
its own; the accumulator has to be given one, and the only choice that keeps
the loop consistent is the element that changes nothing.

It is a convention that pays, and it pays twice over. Program~\ref{prog:F01}
got $a^{0} = 1$ out of the division law and said it was forced rather than
chosen; $a^{0}$ is also a product of no factors, so here is a second and
independent route to the same answer. And Program~\ref{prog:F03}'s law carries
one convention onto the other: the logarithm of the empty product is
$\ln 1 = 0$, which is exactly the empty sum, so $\ln \prod = \sum \ln$ holds
at zero terms with no special case bolted on.
\end{trapbox}

Here is where it stops being tidy-mindedness. A batch arrives in which every
token has been masked out, and your loss line is \code{total / count}.

What is \code{total}, what is \code{count}, and what does that line do?
\nextframe
\end{fr}

\begin{fr}
\begin{ansblock}
\code{total} is $0$, which is the empty sum and is perfectly well defined.
\code{count} is $0$ as well. The line divides zero by zero.
\end{ansblock}

\begin{warning}
\textbf{Your tooling disagrees with itself here, and it disagrees silently.}
Both halves of this were run rather than remembered.

In plain Python, \code{0.0 / 0.0} raises \code{ZeroDivisionError} and the
program stops at the line that caused it. In numpy, dividing one float scalar
by another returns \code{nan}, warns, and carries on; and
\code{np.array([]).mean()} likewise returns \code{nan}, warns, and carries on.

The \emph{invariant} is that one library stops and the other hands you a value
and continues. The exact wording of a numpy warning is a property of the
version you happen to have installed, so it is not quoted here \dash{} and a
warning is not an exception: nothing is waiting for you to read it.
\end{warning}

\begin{aibox}
The asymmetry is the useful part. \textbf{An empty sum is safe and an empty
mean is not}, because the sum is defined and the division is not, so a guard
belongs on the denominator and never on the accumulator.

A \code{nan} loss from one fully masked micro-batch does not stay in that
micro-batch: it reaches the gradients, the gradients reach the weights, and
what you see some minutes later is a model whose every output is \code{nan}
with nothing pointing back at the batch that did it.
\end{aibox}
\end{fr}

% ============================================================
\section{Every loss is a sigma and a divisor}
\label{sec:F04-loss}
\index{loss!as a mean over a batch}
\index{sigma notation!in a loss function}

\begin{fr}
Program~\ref{prog:F03} wrote that sum with a sigma it had never defined, and
said the rest of the loss in words: take the logarithm of the probability the
model gave each token, add them up, divide by the number of tokens, and flip
the sign. At $\val{f03.seq.tokens}$ tokens and $\val{f02.loss.nats}$ nats apiece,
that is the loss on the model card.

Here is the same sentence with every symbol named:
\[ L = -\frac{1}{n}\sum_{i=1}^{n} \ln p_i \]
The sigma is the adding up. The $n$ underneath it is the count of the things
that were added. The minus sign is there because the logarithm of a
probability is negative and a loss is reported positive. Nothing new has
happened: the notation has caught up with what you were already doing.

The shape it lives in is more general than that. Take four predictions and the
four values they were meant to hit:
\[ y = (3,\ -\num{0.5},\ 2,\ 7) \qquad
   \hat{y} = (\num{2.5},\ 0,\ 2,\ 8) \]
Mean squared error is
\[ \mathrm{MSE} = \frac{1}{n}\sum_{i=1}^{n} (\hat{y}_i - y_i)^{2} \]
Work it out.
\nextframe
\end{fr}

\begin{fr}
\begin{ansblock}
The residuals are $-\num{0.5}$, $\num{0.5}$, $0$, $1$; the squares are
$\num{0.25}$, $\num{0.25}$, $0$, $1$; the total is $\val{f04.mse.sum}$ and
over four examples that is $\val{f04.mse}$.
\end{ansblock}

Change the body and leave the rest alone and you have a different loss. Mean
absolute error puts the size of the residual in the body instead of its
square:
\[ \mathrm{MAE} = \frac{1}{n}\sum_{i=1}^{n} \lvert \hat{y}_i - y_i \rvert \]
and on the same four points it comes to $\val{f04.mae}$. The sigma did not
move. The divisor did not move. One term moved.

\yourturn
Write cross-entropy over a batch of $n$ examples in that same shape, given
that the per-example term is the negative logarithm of the probability the
model put on the right answer.
\dotline
\nextframe
\end{fr}

\begin{fr}
\begin{ansblock}
\[ L = -\frac{1}{n}\sum_{i=1}^{n} \ln p_i \]
which is the line this section opened with, reached from the other end.
\end{ansblock}

Three losses side by side. Only one of the three columns changes with the
loss \dash{} the body:

\begin{center}
\begin{tabular}{lll}
\toprule
Loss & The shape & The body \\
\midrule
Mean squared error & $\frac{1}{n}\sum$ & $(\hat{y}_i - y_i)^{2}$ \\
Mean absolute error & $\frac{1}{n}\sum$ & $\lvert \hat{y}_i - y_i \rvert$ \\
Cross-entropy & $\frac{1}{n}\sum$ & $-\ln p_i$ \\
\bottomrule
\end{tabular}
\end{center}

Now use that third body. A batch of four examples, in which
the model put probability $\num{0.9}$, $\num{0.5}$, $\num{0.2}$ and
$\num{0.05}$ on the right answer. What is the loss?
\dotline
\nextframe
\end{fr}

\begin{fr}
\ans{$\val{f04.ce.mean}$}
The four terms come to $\val{f04.ce.sum}$, and there are four of them.

Program~\ref{prog:F03} turned a mean loss into a perplexity by exponentiating
it. Doing that here is worth the detour, because it produces two numbers where
you might have expected one. The exponential of the mean loss is
$\val{f04.ce.ppl}$. The mean of the four per-example perplexities \dash{} that
is, the mean of the four values of $1/p_i$ \dash{} is
$\val{f04.ce.ppl.naive}$.

Which of those two is the number a leaderboard should print, and why?
\nextframe
\end{fr}

\begin{fr}
\begin{ansblock}
The exponential of the mean, $\val{f04.ce.ppl}$. Program~\ref{prog:F03}
settled it: perplexity is $e$ raised to the mean negative log-likelihood, and
an average of perplexities is a different quantity that gets reported in the
same column.
\end{ansblock}

The general statement behind that \dash{} what happens when you average a
curved function of a quantity rather than the quantity itself \dash{} is
Jensen's inequality, and it belongs to Program~\ref{prog:P19}. What this
program owes you is the shape underneath it:

\textbf{every loss here is $\frac{1}{n}\sum$ over a per-example term, and the
only thing that moves from one loss to the next is the term.}

\begin{aibox}
That sentence is a line of the training step. A language model's
cross-entropy is summed over the positions whose label is not the ignore
index, and the divisor is the count of \emph{those} positions \dash{} not the
size of the tensor, which also contains the padding.

Two counts, one of which is larger, and the loss is the sum divided by
whichever one the code picked up. \S\ref{sec:F04-averages} is about what
happens when the divisor is computed in the wrong place.
\end{aibox}
\end{fr}

\begin{fr}
One thing a sigma does not promise.

In exact arithmetic a sum has no order. $a + b + c$ is $c + b + a$, and there
is nothing to discuss. In floating point it has an order, because every
addition rounds: a very small number added to a much larger running total can
round away to nothing whatever, and it does so without an error, a warning or
a mark of any kind. Add a million small gradients into one accumulator and
some of them will not be in the answer.

\begin{warning}
\textbf{This frame quotes no number, deliberately.} How much is lost depends
on the format, the order, the hardware, and which reduction the library chose,
so a figure printed here would be a figure about one machine on one day.

Program~\ref{prog:P01} says why a sum of gradients is not associative and why
two accelerators therefore disagree in the last bits.
Program~\ref{prog:P02} says what summing a million small numbers in the wrong
order costs, and what to do about it.

The sigma is still the right notation. It is a statement about the value; your
machine computes an approximation to it, in an order you did not choose.
\end{warning}
\end{fr}

% ============================================================
\section{When you may not average the averages}
\label{sec:F04-averages}
\index{measurement!micro and macro averages}
\index{troubleshooting!average of averages}

\begin{fr}
No working on this one either. The first number that comes to you.

Two micro-batches have gone through. The first held $1000$ tokens and came
back with a mean loss of $\num{2.0}$. The second held $10$ tokens and came
back with a mean loss of $\num{8.0}$.

\textbf{What is the mean loss over the $1010$ tokens?}
\dotline
\nextframe
\end{fr}

\begin{fr}
\ans{$\val{f04.batch.pooled}$}
Not $\num{5.0}$. Every token contributes its own loss, and there are a great
many more of them in one batch than in the other:
\[ \frac{1000 \times \num{2.0} + 10 \times \num{8.0}}{1010}
   = \val{f04.batch.pooled} \]
The answer that comes first is larger than the true one by a factor of
$\val{f04.batch.factor}$.

\begin{trapbox}
\textbf{Two numbers were averaged that were already averages.} Adding
$\num{2.0}$ to $\num{8.0}$ and halving gives each of them the weight one half
\dash{} which would be right if they stood for equal amounts of evidence, and
one of them stands for a hundred times as much as the other.

Both quantities are real and both have names. The \emph{micro} average pools
the items and divides once, so every token counts the same. The \emph{macro}
average averages the group numbers, so every group counts the same. There is
nothing wrong with either of them.

So the fault is not that you computed the wrong one. The fault is not knowing
which one you are holding, and the reason that happens is that they are
reported in the same column, under the same word, to the same number of
decimal places.
\end{trapbox}
\end{fr}

\begin{fr}
This is not a one-off about batch sizes. The second instance is in every
evaluation harness you have ever run.

An evaluation set is split across two shards. The first holds $100$ items and
the model gets $90$ of them right. The second holds $10$ items and the model
gets $1$ right. Your harness prints the two per-shard accuracies:
$\num{0.90}$ and $\num{0.10}$.

What is the accuracy on the $110$ items?
\nextframe
\end{fr}

\begin{fr}
\begin{ansblock}
\[ \frac{91}{110} = \val{f04.acc.micro} \]
Every item the model got right, over every item there was. Not $\num{0.50}$.
\end{ansblock}

A third instance, in a different unit, where the gap is widest. One run does
$100$ tokens in $1$ second, another $100$ tokens in $100$ seconds: rates of
$100$ and $1$ token per second, whose average is $\num{50.5}$. The rate over
the whole of the work is
\[ \frac{100 + 100}{1 + 100} = \val{f04.rate.micro} \]
tokens per second, smaller by a factor of $\val{f04.rate.factor}$. That figure
is exact and checkable: with equal numerators the pooled rate is the
\emph{harmonic} mean of the two rates, $\val{f04.rate.micro}$ to the digit.

Now say in one sentence what the three instances have in common.
\dotline
\nextframe
\end{fr}

\begin{fr}
\begin{ansblock}
In each of them a mean of ratios was put where a ratio of sums belonged.
\end{ansblock}

Written once and in general, for a group $i$ with numerator $a_i$ and
denominator $b_i$:
\[ \frac{1}{k}\sum_{i=1}^{k} \frac{a_i}{b_i} \ne
   \frac{\sum_{i=1}^{k} a_i}{\sum_{i=1}^{k} b_i} \]
The left is the \emph{macro} average: every group carries weight $1/k$,
whatever it stands for. The right is the \emph{micro} average: every item
carries the same weight, so a group's influence is its size. The two agree
exactly when every $b_i$ is equal, and not otherwise.

Keep this apart from Program~\ref{prog:F03}'s perplexity failure, which looks
like the same thing and is not. There, a curved function was applied to a mean
and the curvature did the damage. Here every step is linear and the damage is
done by the weights. Same smell, different mechanism, and the two do not have
the same fix.

\yourturn
Your harness reports an accuracy for each benchmark and averages the numbers.
Under what condition is that average the accuracy on the pooled items?
\dotline
\nextframe
\end{fr}

\begin{fr}
\begin{ansblock}
When the benchmarks all hold the same number of items. Then equal weights are
the right weights, and the two averages coincide.
\end{ansblock}

\begin{aibox}
Two places this is live, and neither is exotic.

The first is a multi-benchmark leaderboard. Its headline number is a macro
average over benchmarks whose sizes differ by orders of magnitude, so it
answers a different question from the pooled score, and a model can be moved
up it by doing well on the small ones. Program~\ref{prog:P27} asks whether the
gap between two such numbers is larger than the noise in either, and
Program~\ref{prog:P34} makes that the standard a reported number has to meet.

The second is gradient accumulation. A language model's cross-entropy is
normalised by the count of positions that were not masked out, and
micro-batches do not all hold the same number of those. Take the mean within
each micro-batch and then average across the micro-batches, and every
micro-batch gets equal weight however many real tokens it carried: what is
being optimised is not what the equivalent large batch would have optimised.
The denominator belongs to the whole accumulated batch. That one is
Program~\ref{prog:P21}'s.
\end{aibox}
\end{fr}

\mermaidfig{f04-two-averages}{%
  One split of one evaluation set, and two ways to divide. The denominator is
  the whole difference between them: the number of groups, or the number of
  items.}{%
  Macro and micro averaging over two unequal shards.}

% ============================================================
\section{Sequences, recurrences and the geometric series}
\label{sec:F04-sequences}
\index{sequence!recurrence}
\index{sequence!geometric series}

\begin{fr}
A \emph{sequence} is a list with an index: $a_1$, $a_2$, $a_3$, and onwards. A
\emph{recurrence} defines each term from the one before it, which is a loop
that carries a value rather than accumulating one.

You own two of them already, without the words. An arithmetic sequence adds a
fixed amount at each step; a geometric sequence multiplies by a fixed factor:

\begin{center}
\begin{tabular}{lll}
\toprule
 & Recurrence & General formula \\
\midrule
Arithmetic & $a_{n+1} = a_n + d$ & $a_n = a_1 + (n-1)d$ \\
Geometric & $a_{n+1} = r\,a_n$ & $a_n = a_1 r^{\,n-1}$ \\
\bottomrule
\end{tabular}
\end{center}

The recurrence says how to take one step. The general formula says where you
end up without taking any, which is what you want when the number of steps is
large.

Take $a_1 = 3$ and $a_{n+1} = 2 a_n$. What is $a_5$, and what is the general
formula?
\nextframe
\end{fr}

\begin{fr}
\begin{ansblock}
The terms are $3$, $6$, $12$, $24$, $48$, so $a_5 = 48$; and
$a_n = 3 \times 2^{\,n-1}$, which delivers $3 \times 2^{4} = 48$ without
walking the sequence at all.
\end{ansblock}

Now set Program~\ref{prog:F03} beside it. Take logarithms of a geometric
sequence:
\[ \ln\!\left(a_1 r^{\,n-1}\right) = \ln a_1 + (n-1)\ln r \]
and what is on the right is an arithmetic sequence in $n$, with the fixed
amount $\ln r$ added at every step. \textbf{A logarithm turns a geometric
sequence into an arithmetic one} \dash{} which is Program~\ref{prog:F03}'s
logarithmic axis, reached from the other end.

A loss halves every thousand steps. What shape is that curve on an axis whose
ticks read $1$, $\num{0.1}$, $\num{0.01}$?
\nextframe
\end{fr}

\begin{fr}
\ans{A straight line}
Every thousand steps multiplies by the same factor, so it moves the same
distance down the axis. A constant ratio is a constant distance, which is the
whole of what a logarithmic axis does.

Now sum a geometric sequence. This is a closed form the program wants you to be
able to produce rather than recall, and producing it takes three
lines. Write $S$ for the sum of the first $n$ terms of $1$, $r$, $r^{2}$, and
so on:
\[ S = \sum_{k=0}^{n-1} r^{k} = 1 + r + r^{2} + \cdots + r^{\,n-1} \]
Multiply the whole line by $r$ and write it underneath, shifted along by one:
\[ rS = r + r^{2} + \cdots + r^{\,n-1} + r^{\,n} \]
Subtract the second from the first. Everything in the middle cancels, which is
the move Program~\ref{prog:F02} calls collecting like terms:
\[ S - rS = 1 - r^{\,n} \qquad\text{and therefore}\qquad S = \blank \]

Then evaluate $\sum_{k=0}^{9} (\num{0.9})^{k}$.
\dotline
\nextframe
\end{fr}

\begin{fr}
\begin{ansblock}
\[ S = \frac{1 - r^{\,n}}{1 - r} \]
At $r = \num{0.9}$ and $n = 10$ that is $\val{f04.geom.ten}$.
\end{ansblock}

Now let the number of terms grow. If $\lvert r \rvert < 1$ then $r^{\,n}$
falls towards nothing, the numerator settles at $1$, and the sum settles too:
\[ \sum_{k=0}^{\infty} r^{k} = \frac{1}{1-r} \]
which at $r = \num{0.9}$ is $\val{f04.geom.inf}$: ten terms have already
brought in $\val{f04.geom.ten}$ of it, and every remaining term put together
supplies the rest.

\begin{rigourbox}
\textbf{That the sum settles at all is not proved here.} A sum of infinitely
many terms is defined as the limit of its partial sums, the limit exists
exactly when $\lvert r \rvert < 1$, and the proof of that is a page of a first
analysis course.

What this book needs from it is the arithmetic and the boundary. Below $1$ in
size, the total is finite and the formula holds. At $1$ or above, the terms do
not shrink and there is no total to have.
\end{rigourbox}
\end{fr}

\begin{fr}
The other direction is a product rather than a sum, and it decides whether a
deep stack works at all.

Take a factor applied once per layer, over fifty layers:
\[ (\num{0.9})^{50} = \val{f04.decay.fifty} \qquad
   (\num{1.1})^{50} = \val{f04.grow.fifty} \]
Two factors $\num{0.2}$ apart at the start, about $\val{f04.decay.ratio}$
apart at the far end, and neither of them extreme.

\yourturn
Compute $50 \ln \num{0.9}$, and say what it tells you about the first of those
two numbers.
\dotline
\nextframe
\end{fr}

\begin{fr}
\begin{ansblock}
$50 \ln \num{0.9} = \val{f04.decay.log}$, so the product is
$e^{\val{f04.decay.log}}$, which is $\val{f04.decay.fifty}$. Fifty
multiplications became one multiplication.
\end{ansblock}

That step is the identity joining this program's two symbols, and it is
Program~\ref{prog:F03}'s first law written over a pi rather than over two
factors:
\[ \ln \prod_{i=1}^{n} x_i = \sum_{i=1}^{n} \ln x_i \]
A product of a thousand probabilities is computed as a sum of a thousand
logarithms for exactly this reason, and Program~\ref{prog:F03} measured what
becomes of anybody who does it the other way.

A signal crossing fifty layers is multiplied once per layer, so a per-layer
factor is not added up over the stack \dash{} it is raised to the fiftieth
power, and that is why $\num{0.9}$ and $\num{1.1}$ per layer are not
near-misses of one another. Why the factors multiply rather than add is
Program~\ref{prog:F12}'s: a composition of layers is differentiated by the
chain rule, and a chain rule is a product.
\end{fr}

% ============================================================
\section{The exponential moving average}
\label{sec:F04-ema}
\index{moving average!exponential}
\index{moving average!bias correction}

\begin{fr}
Here is a problem you have met on every dashboard you have ever looked at. The
loss curve jumps about from step to step and what you want to see is the trend
underneath it. You could keep the last hundred values and average them, but
that means keeping a hundred values, at every step, for every quantity you are
watching.

There is a two-line alternative that keeps nothing at all:

\begin{python}
m = 0.0
for g in stream:
    m = beta * m + (1 - beta) * g
\end{python}

One number of state and one multiply-add per observation. It is called an
\emph{exponential moving average}, and this section takes it apart because it
is the whole of momentum and half of Adam.

Take $\beta = \num{0.9}$. Quickly, and with no working: $m$ currently reads
$\num{1.0}$ and the new $g$ is $\num{11.0}$. What does $m$ read after one
step?
\dotline
\nextframe
\end{fr}

\begin{fr}
\ans{$\num{2.0}$}
\[ \num{0.9} \times \num{1.0} + \num{0.1} \times \num{11.0}
   = \num{0.9} + \num{1.1} = \num{2.0} \]

\begin{trapbox}
\textbf{The other answer is $\num{10.0}$, and it comes from reading $\beta$ as
how much of the new value gets in.} It is the other way round.
\textbf{$\beta$ weights the old value}, so $\beta = \num{0.9}$ keeps
nine-tenths of everything that has happened so far and lets one-tenth of the
news in.

What makes the misreading so easy is that $\beta$ is the number written in the
configuration file and the coefficient on $g$ is written nowhere: it is
$1 - \beta$, and it is computed. The number you set is the one least likely to
feel as though it belongs to the past.

Everything about the average inverts with it. Raising $\beta$ makes the
average \emph{slower} and smoother, not quicker to react, because more of the
old value survives each step.
\end{trapbox}

Which forgets a change faster, $\beta = \num{0.9}$ or $\beta = \num{0.999}$?
\nextframe
\end{fr}

\begin{fr}
\ans{$\beta = \num{0.9}$}
It lets a tenth of the news in at each step, where $\beta = \num{0.999}$ lets
in a thousandth.

Now unroll the recurrence into itself, which is the frame this whole program
has been walking towards. Substitute the definition of $m_{t-1}$ into the
definition of $m_t$, then do it again, and again:
\[ m_t = (1-\beta)\,g_t + (1-\beta)\beta\,g_{t-1}
       + (1-\beta)\beta^{2} g_{t-2} + \cdots + \beta^{\,t} m_0 \]
which collects into
\[ m_t = (1-\beta)\sum_{k=0}^{t-1} \beta^{k} g_{t-k} + \beta^{\,t} m_0 \]

\textbf{The two-line loop is a weighted sum of every value it has ever seen.}
The weights are \S\ref{sec:F04-sequences}'s geometric sequence, and it is
\S\ref{sec:F04-sequences}'s series that says what they come to. Nothing new
was needed to get here, which is why the sequences came first.

At $\beta = \num{0.9}$, what weight sits on the newest value $g_t$, and what
weight sits on the one before it?
\[ \text{newest: } \blank \qquad\qquad \text{the one before: } \blank \]
\dotline
\nextframe
\end{fr}

\begin{fr}
\begin{ansblock}
$1 - \beta = \num{0.10}$ on the newest, and $(1-\beta)\beta = \num{0.09}$ on
the one before it.
\end{ansblock}

Which lets the folklore be checked instead of repeated.
\enquote{$\beta = \num{0.9}$ averages the last ten} is what everybody says.
The weights are $(1-\beta)\beta^{k}$, so the last ten carry
\[ 1 - (\num{0.9})^{10} = \val{f04.ema.w10} \]
of the total weight, the last twenty carry $\val{f04.ema.w20}$, and the rest
is spread out behind them without end.

Three different numbers answer \emph{how far back does this thing look}, and
none of them is a count of ten observations:

\begin{center}
\begin{tabular}{lll}
\toprule
Reading & At $\beta = \num{0.9}$ & \\
\midrule
Effective window, $1/(1-\beta)$ & $\val{f04.geom.inf}$ & a scale, not a count \\
Centre of mass, $\beta/(1-\beta)$ & $\val{f04.ema.com}$ & steps back \\
Half-life, $\ln \frac{1}{2} / \ln \beta$ & $\val{f04.ema.halflife}$ & steps \\
\bottomrule
\end{tabular}
\end{center}

The half-life is the one to carry, because it is the smallest of the three and
it answers the question you actually have: \textbf{half the weight is in the
last $\val{f04.ema.halflife.steps}$ observations.} The ten of the folklore is
the first row, which is the scale the weights are built on rather than a count
of anything.

And $\val{f04.ema.w10}$ is not a second sighting of
\S\ref{sec:F04-sequences}'s $\val{f04.geom.ten}$ by coincidence. It is that
sum multiplied by $1 - \beta$, which is what the weights are.
\end{fr}

\mermaidfig{f04-ema-weights}{%
  Where an exponential moving average keeps its weight. The values are not
  forgotten at some cut-off; each is worth $\beta$ of the one after it, for
  ever.}{%
  The weights of an exponential moving average, and its half-life.}

\begin{fr}
\yourturn
The initialisation, which the unrolling has a term for and nobody ever looks
at.

$m$ starts at $0$, and every observation that arrives is exactly $1$. What
does $m$ read after one step at $\beta = \num{0.9}$, and what does it read
after one step at $\beta = \num{0.999}$?
\dotline
\nextframe
\end{fr}

\begin{fr}
\begin{ansblock}
$\num{0.1}$ and $\num{0.001}$.
\end{ansblock}

The average of a constant $1$ reads a tenth of the truth after one step, and a
thousandth of it at $\beta = \num{0.999}$. That is not a slow start; it is a
wrong number, and it stays wrong for a while. After $t$ steps the reading is
$1 - \beta^{\,t}$ of the truth \dash{} the same $1 - \beta^{\,t}$ that came
out of the geometric sum, because it is the fraction of the weight that has
arrived so far.

Which hands you the fix, since you know what fraction has arrived. Divide by
it:
\[ \hat{m}_t = \frac{m_t}{1 - \beta^{\,t}} \]
On a constant sequence that is exactly right at every $t$ in the algebra:
$m_t$ is $1 - \beta^{\,t}$ times the constant, and the division puts it back.
In binary64 it is right to about the last place, which is the ordinary
arithmetic of floats and not a defect in the correction.

\begin{aibox}
This is Adam's bias correction, and the two coefficients are Adam's two.
Kingma and Ba's paper sets $\beta_1 = \num{0.9}$ for the first moment and
$\beta_2 = \num{0.999}$ for the second, and derives this same division from
the initialisation at zero.

Their half-lives are $\val{f04.ema.halflife.steps}$ steps and
$\val{f04.ema.b999.steps}$ steps, so at any early step the two moments are
biased by very different amounts, and each is corrected by its own
$1 - \beta^{\,t}$. It is not a detail that matters for the first three steps
and then stops: at $\beta = \num{0.999}$ the correction is doing real work for
hundreds of them.
\end{aibox}

Your smoothed loss curve starts well below the real loss and climbs to meet it
over the first few hundred steps. What is happening?
\nextframe
\end{fr}

\begin{fr}
\begin{ansblock}
The bias, not the model. The average is still climbing out of its zero
initialisation, and $1 - \beta^{\,t}$ says how far out of it it has got.
Divide by that and the climb is gone.
\end{ansblock}

Now cash the program in. \textbf{Momentum is an exponential moving average of
the gradient, stepped along in place of the gradient itself}, so the optimiser
moves along a smoothed direction rather than along the last noisy one.
\textbf{Adam is this
recurrence run twice} \dash{} once on the gradient and once on its square
\dash{} with the correction of the previous frame on each.

None of that used calculus. Every line of this section has been arithmetic on
a list of numbers, and it is yours now rather than in Part~V.
Program~\ref{prog:P20} supplies the gradient and the rest of the optimiser;
what it does not have to supply is the average.

\begin{warning}
\textbf{Two things about that coefficient are convention rather than
mathematics, and the second one changes the length of the step.} The first you
have met: it may weight the old value, as it does here, or the new one, and
$\num{0.9}$ under one reading is $\num{0.1}$ under the other.

The second is the $1 - \beta$ itself, which is often not written at all.
Keep $v \leftarrow \beta v + g$ and every weight in the unrolled sum has been
multiplied by $\frac{1}{1-\beta}$: the direction is identical and the vector
is ten times as long at $\beta = \num{0.9}$, a thousand times at
$\beta = \num{0.999}$. A constant factor on a step is indistinguishable from
a change of step size, which is why both forms are in use and neither is
wrong. What does not survive the change is a step size carried across from the
other one. Program~\ref{prog:P20} is where that has consequences.

Nothing about the name tells you which you have. Read the line that updates
the state: which side the coefficient sits on, and whether the other side
carries a factor at all.
\end{warning}

Somebody changes a running average from $\beta = \num{0.99}$ to
$\beta = \num{0.9}$ in order to \enquote{smooth it more}. What happens?
\nextframe
\end{fr}

\begin{fr}
\begin{ansblock}
It smooths \emph{less}. The half-life falls from
$\val{f04.ema.b99.halflife}$ steps to $\val{f04.ema.halflife}$, so the curve
looks back about a tenth as far and becomes correspondingly noisier.
\end{ansblock}

\medskip
\noindent
\textbf{What comes next.} Program~\ref{prog:F05} takes the function itself as
its subject, and a sequence is a function whose input is an index.
Program~\ref{prog:F13} is the continuous twin of this program's sigma: an
integral accumulates over a range where a sum accumulates over a list.
Program~\ref{prog:P03} prices a model in sigma notation and needs nothing more
than the arithmetic you have. Program~\ref{prog:P12} counts, which is a pi
wearing a different hat. Program~\ref{prog:P19} has the general statement
behind \S\ref{sec:F04-loss}'s perplexity, and Program~\ref{prog:P20} attaches
a gradient to \S\ref{sec:F04-ema}'s recurrence. And
Program~\ref{prog:P25} answers the question this program has carefully not
asked: how accurate is a mean over $n$ things, and how much does another
thousand of them buy you?
\end{fr}

% ============================================================
\begin{summarybox}
\sumitem{1--2}{\result{\textbf{A sigma is a loop written short.} Its four parts are the
  index, where the index starts, where it stops, and the body \dash{} and the
  body is an expression in the index, worked once for each value the index
  takes}.}
\sumitem{3--4}{\result{$\sum_{i=1}^{n} w_i x_i$ is a \emph{weighted sum}: one unit of a
  linear layer, and the shape Program~\ref{prog:P05} is built on. A sigma
  splits over a sum, and a constant factor comes out of the front. Both are the
  loop rearranged}.}
\sumitem{5--6}{\result{\textbf{$\sum a_i b_i$ is not $(\sum a_i)(\sum b_i)$}, and at
  $a = (1, 2)$, $b = (3, 4)$ the two are $11$ and $21$. It is a rule about a
  sum inside the sigma being carried across a product, which is the reasoning
  Programs~\ref{prog:F01}, \ref{prog:F02} and~\ref{prog:F03} each caught under
  a different sign}.}
\sumitem{7--8}{\result{The number of terms is $\text{stop} - \text{start} + 1$, because
  both limits are included. $\sum_{i=0}^{n}$ has $n+1$ terms and
  $\sum_{i=0}^{n-1}$ has $n$, which is what \code{range(n)} walks}.}
\sumitem{9}{\result{Re-indexing moves the limits one way and the subscript the other:
  $\sum_{i=1}^{n} x_i = \sum_{j=0}^{n-1} x_{j+1}$. Note that English calls both
  the raised number and the lowered one the \emph{index}, and Polish does not}.}
\sumitem{10--11}{\result{A causal mask computes $\frac{n(n+1)}{2}$ query--key pairs
  \dash{} $\val{f04.causal.pairs}$ at $n = 2048$ against $\val{f04.causal.sq}$
  for the full square. The ratio $\frac{n+1}{2n}$ falls towards a half from
  above and never reaches it: $\val{f04.causal.ratio.eight}$ at $n = 8$}.}
\sumitem{12--13}{\result{$\prod$ is the same machine with multiplication in it, and a
  factorial is a pi with the index in the body: $n! = \prod_{i=1}^{n} i$}.}
\sumitem{14--15}{\result{\textbf{The empty sum is $0$ and the empty product is $1$},
  because each is the identity element of its operation and no other value
  keeps the loop consistent. $a^{0} = 1$ is the empty product, which
  corroborates Program~\ref{prog:F01}'s division-law argument by a second
  route, and $\ln 1 = 0$ carries one convention onto the other}.}
\sumitem{16}{\result{An empty \emph{sum} is safe and an empty \emph{mean} is not.
  \code{0.0 / 0.0} raises in plain Python; numpy returns \code{nan}, warns, and
  carries on. The guard belongs on the denominator}.}
\sumitem{17--19}{\result{\textbf{Every loss here is $\frac{1}{n}\sum$ over a
  per-example term}, and only the term moves: a squared residual, its size, or
  $-\ln p_i$. On four points, mean squared error is $\val{f04.mse}$ and mean
  absolute error is $\val{f04.mae}$}.}
\sumitem{20--21}{\result{Perplexity is the exponential of the mean loss,
  $\val{f04.ce.ppl}$ on that batch. The mean of the per-example perplexities is
  $\val{f04.ce.ppl.naive}$, which is a different quantity;
  Program~\ref{prog:P19} has the inequality behind it}.}
\sumitem{22}{\result{In exact arithmetic a sum has no order and in floating point it
  has one.} This program prints no number for it, because the number would be
  about one machine; Programs~\ref{prog:P01} and~\ref{prog:P02} own it.}
\sumitem{23--26}{\result{\textbf{A mean of ratios is not the ratio of the sums.} Two
  batches at $\num{2.0}$ and $\num{8.0}$ pool to $\val{f04.batch.pooled}$ and
  not $\num{5.0}$; two shards at $\num{0.90}$ and $\num{0.10}$ pool to
  $\val{f04.acc.micro}$ and not $\num{0.50}$; two runs at $100$ and $1$ token
  per second pool to $\val{f04.rate.micro}$ and not $\num{50.5}$}.}
\sumitem{27--28}{\result{The macro average weights groups equally, the micro average
  weights items equally, and they agree exactly when every denominator is
  equal. Both are legitimate; the fault is reporting one as the other}.}
\sumitem{29--32}{\result{A recurrence takes one step and a general formula skips them
  all. $\sum_{k=0}^{n-1} r^{k} = \frac{1-r^{n}}{1-r}$, which is
  $\val{f04.geom.ten}$ at $r = \num{0.9}$ and $n = 10$ and tends to
  $\frac{1}{1-r} = \val{f04.geom.inf}$}.}
\sumitem{33--34}{\result{$\ln \prod_i x_i = \sum_i \ln x_i$, which is
  Program~\ref{prog:F03}'s first law written over a pi. Fifty layers at
  $\num{0.9}$ give $\val{f04.decay.fifty}$ and at $\num{1.1}$ give
  $\val{f04.grow.fifty}$: two factors $\num{0.2}$ apart are about
  $\val{f04.decay.ratio}$ apart at the far end}.}
\sumitem{35--38}{\result{An exponential moving average is
  $m \leftarrow \beta m + (1-\beta) g$, and unrolled it is a weighted sum of
  everything it has seen, with weights $(1-\beta)\beta^{k}$.
  \textbf{$\beta$ weights the old value.} At $\beta = \num{0.9}$ the last ten
  carry $\val{f04.ema.w10}$ of the weight and half of it sits in the last
  $\val{f04.ema.halflife.steps}$}.}
\sumitem{39--42}{\result{Starting at $0$ biases it: after $t$ steps it reads
  $1 - \beta^{\,t}$ of the truth, and dividing by that removes the bias
  exactly. \textbf{That is momentum, and with the recurrence run twice it is
  Adam}, and none of it needed calculus}.}
\end{summarybox}

\canyou

% ============================================================
\begin{testexercises}
\item Write $\sum_{i=1}^{5} (2i - 1)$ out in full and evaluate it.
  \answerto{$1 + 3 + 5 + 7 + 9 = 25$.}
\item How many terms are there in $\sum_{i=3}^{9} x_i$?
  \answerto{$7$, from $9 - 3 + 1$. Both limits are included.}
\item Rewrite $\sum_{i=1}^{6} x_i$ so that the index starts at $0$.
  \answerto{$\sum_{j=0}^{5} x_{j+1}$.}
\item Write this loop as a sigma:
  \begin{center}
  \code{total = 0.0}\\
  \code{for c, p in zip(counts, prices):}\\
  \code{\quad total = total + c * p}
  \end{center}
  \answerto{$\sum_{i=1}^{n} c_i p_i$, a weighted sum.}
\item Evaluate $\prod_{i=1}^{5} i$.
  \answerto{$120$, which is $5!$.}
\item What do \code{sum([])} and \code{math.prod([])} return, and why those two
  values?
  \answerto{$0$ and $1$. The loop body never runs, so each returns its
  accumulator, and the accumulator has to be the value that leaves the next
  term alone.}
\item Write mean squared error over $n$ examples in sigma notation, and say
  which part of it changes when you switch to mean absolute error.
  \answerto{$\frac{1}{n}\sum_{i=1}^{n}(\hat{y}_i - y_i)^{2}$, and only the body
  changes \dash{} to $\lvert \hat{y}_i - y_i \rvert$. The sigma and the divisor
  stay where they are.}
\item One batch of $400$ tokens has mean loss $\num{1.5}$ and another of $100$
  tokens has mean loss $\num{3.5}$. What is the mean loss over the $500$
  tokens?
  \answerto{$\frac{400 \times \num{1.5} + 100 \times \num{3.5}}{500}
  = \frac{950}{500} = \num{1.9}$, not the $\num{2.5}$ you get by averaging the
  two.}
\item One shard scores $30$ of $50$ and another $30$ of $150$. What is the
  accuracy over all the items, and what is the average of the two shard
  accuracies?
  \answerto{$\frac{60}{200} = \num{0.30}$ pooled, against
  $\frac{\num{0.60} + \num{0.20}}{2} = \num{0.40}$ macro. The larger shard is
  the worse one and the macro average hides it.}
\item Evaluate $\sum_{k=0}^{3} r^{k}$ at $r = \frac{1}{2}$, and say what it
  tends to as the number of terms grows.
  \answerto{$\frac{1 - 1/16}{1/2} = \frac{15}{8} = \num{1.875}$, tending to
  $\frac{1}{1 - 1/2} = 2$.}
\item A running average with $\beta = \num{0.8}$ reads $5$, and the next
  observation is $10$. What does it read after the step?
  \answerto{$\num{0.8} \times 5 + \num{0.2} \times 10 = 4 + 2 = 6$.}
\item An average with $\beta = \num{0.5}$ starts at $0$ and every observation
  is $1$. What does it read after two steps, and what does the bias correction
  give?
  \answerto{$\num{0.5}$ then $\num{0.75}$, which is $1 - \beta^{2}$ of the
  truth. Dividing by $1 - \beta^{2} = \num{0.75}$ gives $1$.}
\end{testexercises}

\begin{furtherproblems}
\item Using the loop rather than a rule, argue that
  $\sum_{i=1}^{n}(a_i + b_i) = \sum_{i=1}^{n} a_i + \sum_{i=1}^{n} b_i$.
  \answerto{Both sides add the same $2n$ numbers; only the order of the
  additions differs, and in exact arithmetic that does not matter. See frame
  $22$ for the one sense in which it does.}
\item Evaluate $\sum_{i=1}^{100} i$ by pairing rather than by adding.
  \answerto{$\frac{100 \times 101}{2} = 5050$. Pair the first with the last,
  the second with the second-last, and every pair comes to $101$.}
\item How many query--key pairs does a causal mask compute at $n = 1024$, and
  what fraction of the full square is that?
  \answerto{$\frac{1024 \times 1025}{2} = 524800$, which is
  $\frac{1025}{2048}$ of the square \dash{} just over a half, and above it.}
\item Show that $0! = 1$ twice: once from the empty product, and once from
  $n! = n \times (n-1)!$ read downwards from $1! = 1$.
  \answerto{A product of no factors is the identity, $1$. And
  $1! = 1 \times 0!$ forces $0! = 1$ if the recurrence is to hold at $n = 1$.}
\item Your loss line is \code{total / count} and a batch can arrive fully
  masked. Write the guard, and say why it belongs on \code{count} rather than
  on \code{total}.
  \answerto{Skip the batch, or return a zero loss, when \code{count} is zero.
  \code{total} is a perfectly well-defined empty sum; it is the division that
  has no answer.}
\item Take $y = (0, 0, 0, 10)$ and $\hat{y} = (0, 0, 0, 0)$. Compute mean
  squared error and mean absolute error, and say which is moved more by the
  single outlier.
  \answerto{$\frac{100}{4} = 25$ and $\frac{10}{4} = \num{2.5}$. Squaring the
  residual before averaging makes one large error dominate the sum, which is
  the choice you are making when you pick the loss.}
\item Three benchmarks of $1000$, $100$ and $10$ items give scores
  $\num{0.5}$, $\num{0.6}$ and $\num{1.0}$. What is the macro average, and what
  is the accuracy over all the items?
  \answerto{Macro is $\frac{\num{2.1}}{3} = \num{0.70}$; pooled is
  $\frac{570}{1110}$, about $\num{0.51}$. The two differ by more than most
  reported gaps between models.}
\item Three runs each do $100$ tokens, taking $1$, $2$ and $4$ seconds. Give
  the mean of the three rates, the rate over the whole of the work, and the
  harmonic mean of the three rates.
  \answerto{Rates $100$, $50$, $25$; their mean is about $\num{58.3}$; the
  pooled rate is $\frac{300}{7}$, about $\num{42.9}$; and the harmonic mean is
  $\frac{3}{\num{0.07}}$, the same $\num{42.9}$, because the numerators are
  equal.}
\item Evaluate $\sum_{k=0}^{4} 2^{k}$ with the same formula, and say why it
  still works when $r > 1$ but the infinite sum does not exist.
  \answerto{$\frac{1 - 32}{1 - 2} = 31$. The derivation never assumed anything
  about $r$; the \emph{limit} needs $r^{n}$ to fall away, and at $r = 2$ it
  grows.}
\item Why is the probability of a fifty-token sequence computed as a sum of
  logarithms rather than as a product?
  \answerto{Because $\ln \prod x_i = \sum \ln x_i$ and the sum stays inside
  the range a float holds. Program~\ref{prog:F03} measured the length at which
  the product reaches exactly zero.}
\item Show that the weights $(1-\beta)\beta^{k}$ add to $1$ as the number of
  terms grows.
  \answerto{$(1-\beta)\sum_{k=0}^{\infty}\beta^{k}
  = (1-\beta) \times \frac{1}{1-\beta} = 1$. The average is an average and not
  merely a sum.}
\item What does an exponential moving average do at $\beta = 0$, and at
  $\beta = 1$?
  \answerto{At $\beta = 0$ it is the latest observation and has no memory. At
  $\beta = 1$ it never moves off its initial value and has no news. Everything
  useful is strictly between.}
\item What does the bias correction $\frac{1}{1 - \beta^{\,t}}$ do as $t$
  grows, and what does that mean in practice?
  \answerto{$\beta^{\,t}$ falls away, so the factor tends to $1$ and the
  correction fades out. It matters at the start and costs nothing later.}
\item Over how many observations does half the weight of an average at
  $\beta = \num{0.999}$ lie, and how does that compare with $\beta = \num{0.9}$?
  \answerto{$\val{f04.ema.b999.steps}$ observations, from a half-life of
  $\val{f04.ema.b999.halflife}$, against $\val{f04.ema.halflife.steps}$ at
  $\beta = \num{0.9}$ \dash{} a hundredfold difference from a coefficient that
  moved by less than a tenth.}
\item And at $\beta = \num{0.99}$?
  \answerto{$\val{f04.ema.b99.steps}$ observations, from a half-life of
  $\val{f04.ema.b99.halflife}$. Each nine added to the coefficient multiplies
  the memory by about ten.}
\item Write the two lines of momentum from the recurrence alone, given that the
  observation arriving at each step is the gradient $g_t$.
  \answerto{$m \leftarrow \beta m + (1-\beta) g$, then step along $m$ instead
  of along $g$. Written $v \leftarrow \beta v + g$ it is the same direction and
  $\frac{1}{1-\beta}$ times as long, and the step size absorbs the difference.
  Program~\ref{prog:P20} supplies the gradient and the step size.}
\item Rewrite $\sum_{i=0}^{n-1} x_{i+1}$ so that the index starts at $1$.
  \answerto{$\sum_{j=1}^{n} x_{j}$. The limits went up by one and the
  subscript came down by one.}
\item A run reports $\num{0.5}$ seconds per token on $10$ tokens and
  $\num{0.1}$ seconds per token on $100$ tokens. What is the time per token
  over the $110$?
  \answerto{$\frac{5 + 10}{110} = \frac{15}{110}$, about $\num{0.14}$ seconds
  per token. Averaging $\num{0.5}$ and $\num{0.1}$ gives $\num{0.30}$, which
  is more than twice as large.}
\end{furtherproblems}
